\documentclass[10pt,a4paper]{amsart}
\usepackage[margin=19mm]{geometry}
\usepackage{amsmath,amssymb,mathtools}
\usepackage[scr=rsfs,cal=euler]{mathalfa}
\usepackage{xfrac}
\usepackage[unicode,hidelinks,bookmarks=false]{hyperref}
\newcommand{\ie}{i.e.\ }
\newcommand{\cf}{cf.\ }
\newcommand{\resp}{respectively\ }
\newcommand{\Subsection}{\subsection}
\providecommand{\nobreakdash}{\nobreak\hbox{-}}
\setlength{\emergencystretch}{2em}
\multlinegap0pt
\usepackage{amssymb}
\newtheorem{proposition}{Proposition}
\title{Some applications of Choi polynomials of linear maps}
\author{Minh Toan Ho, Thanh Hieu Le, Cong Trinh Le, Hiroyuki Osaka}
\date{Source: arXiv:2604.27034v1 (CC BY 4.0)}
\begin{document}
\maketitle

\noindent\emph{Source and licence.} Some applications of Choi polynomials of linear maps. Version arXiv:2604.27034v1 (CC BY 4.0), distributed by arXiv under the Creative Commons Attribution 4.0 International licence, http://creativecommons.org/licenses/by/4.0/, as recorded in the arXiv licence metadata for this exact version and shown on the abstract page https://arxiv.org/abs/2604.27034v1. Full licence notice: \texttt{licenses/arxiv-2604.27034-CC-BY-4.0.txt}.\par\smallskip

\noindent\textit{Editorial scope.} Counted unit: the article's proposition proposition (source0.tex lines 1650--1651). The article's complete original proof of it is delivered with it (source0.tex lines 1653--1771, one \texttt{proof} environment of 119 lines). No mathematics is written, repaired or re-derived: the statement and the proof are the article's own lines, in the article's own notation, and no retained line is substituted or re-flowed.  No further source-local context is needed: every cross-reference of every retained line resolves inside the two delivered spans.  Nothing was omitted from any retained span, so the package carries no omission record.  No retained line contains a citation, so the wrapper carries no bibliography and no bibliography entry.  No retained line contains a figure environment, an image-inclusion command or drawing code, so no image is delivered and the package declares no asset dependency.  Editorial formatting only, in addition: an AMS \texttt{amsart} preamble that restates the article's own declarations and macros used by the retained lines, namely the environments \texttt{proposition} on separate counters, together with the standard packages those lines need.  The retained environments print in this document as Proposition 1 in order of appearance, whereas the article prints the same items with the numbering of its own full document.  The complete native source of the article as distributed in the arXiv e-print for this exact version is bundled unchanged as \texttt{sources/199441/source0.tex}.
\begin{proposition} $ P_{\tau_{4,1}}(x,y)\ge 0 $  for all $x,y\in\mathbb C^4$.
\end{proposition}

\begin{proof}
Let
\[
x=(x_1,x_2,x_3,x_4)^T,\qquad y=(y_1,y_2,y_3,y_4)^T,
\]
and set
\[
a_i:=|x_i|^2,\qquad b_i:=|y_i|^2
\qquad (i=1,2,3,4).
\]
Indices are taken cyclically modulo $4$, so $a_0=a_4$.

From the explicit formula for $\tau_{4,1}(xx^*)$, we have
\[
P_{\tau_{4,1}}(x,y)
=
\sum_{i=1}^4 (3a_i+a_{i-1})\,b_i
-\left|\sum_{i=1}^4 \overline{x_i}y_i\right|^2.
\]
Thus it suffices to prove
\[
\left|\sum_{i=1}^4 \overline{x_i}y_i\right|^2
\le
\sum_{i=1}^4 (3a_i+a_{i-1})\,b_i.
\]

Assume first that
\[
3a_i+a_{i-1}>0
\qquad (i=1,2,3,4).
\]
Then
\[
\sum_{i=1}^4 \overline{x_i}y_i
=
\sum_{i=1}^4
\frac{\overline{x_i}}{\sqrt{3a_i+a_{i-1}}}\,
\sqrt{3a_i+a_{i-1}}\,y_i.
\]
By the Cauchy--Schwarz inequality,
\[
\left|\sum_{i=1}^4 \overline{x_i}y_i\right|^2
\le
\left(\sum_{i=1}^4 \frac{a_i}{3a_i+a_{i-1}}\right)
\left(\sum_{i=1}^4 (3a_i+a_{i-1})\,b_i\right).
\]
Therefore it remains to prove that
\[
\sum_{i=1}^4 \frac{a_i}{3a_i+a_{i-1}}\le 1.
\]

Assume now that all $a_i>0$, and define
\(\displaystyle
u_i:=\frac{a_{i-1}}{a_i}>0
\qquad (i=1,2,3,4).
\)
Then
\(
u_1u_2u_3u_4=1,
\)
and
\(
\dfrac{a_i}{3a_i+a_{i-1}}=\dfrac{1}{3+u_i}.
\)
Hence it is enough to show that
\[
\sum_{i=1}^4 \frac{1}{3+u_i}\le 1,
\qquad\text{whenever }u_1u_2u_3u_4=1.
\]
Multiplying both sides by $\prod_{i=1}^4(3+u_i)$, this is equivalent to
\[
\prod_{i=1}^4(3+u_i)-\sum_{i=1}^4 \prod_{j\ne i}(3+u_j)\ge 0.
\]
Since $u_1u_2u_3u_4=1$,  the left-hand side becomes
\[
2\sum_{1\le i<j<k\le 4}u_i u_j u_k
+3\sum_{1\le i<j\le 4}u_i u_j
-26.
\]
Applying the arithmetic--geometric mean inequality, we obtain 
\[
\frac{1}{6}\sum_{1\le i<j\le 4}u_i u_j
\ge
\left(\prod_{1\le i<j\le 4}u_i u_j\right)^{1/6}
=
(u_1u_2u_3u_4)^{1/2}=1,
\]
and 
\[
\frac{1}{4}\sum_{1\le i<j<k\le 4}u_i u_j u_k
\ge
\left(\prod_{1\le i<j<k\le 4}u_i u_j u_k\right)^{1/4}
=
(u_1u_2u_3u_4)^{3/4}=1.
\]
Therefore,
\[
2\sum_{1\le i<j<k\le 4}u_i u_j u_k
+3\sum_{1\le i<j\le 4}u_i u_j
\ge 26,
\]
and hence
\(\displaystyle
\sum_{i=1}^4 \frac{1}{3+u_i}\le 1.
\)
It follows that
\(\displaystyle
\sum_{i=1}^4 \frac{a_i}{3a_i+a_{i-1}}\le 1.
\)
Therefore
\[
P_{\tau_{4,1}}(x,y)
=
\sum_{i=1}^4 (3a_i+a_{i-1})\,b_i
-\left|\sum_{i=1}^4 \overline{x_i}y_i\right|^2
\ge 0.
\]
This proves the claim when all $a_i>0$. If some $a_i=0$, the conclusion follows by continuity. 
\end{proof}

\end{document}
